\documentclass[11pt]{article}
\usepackage[a4paper,margin=2.5cm]{geometry}
\usepackage{amsmath}
\usepackage{amssymb}
\usepackage{enumitem}
\usepackage{parskip}
\setlist[enumerate,2]{label=(\alph*), itemsep=1pt, topsep=2pt}
\setlength{\parindent}{0pt}
\newcommand{\correct}[1]{\textbf{#1}\;$\checkmark$}
\begin{document}
\section*{Nash equilibrium}
\begin{enumerate}[itemsep=6pt]
  \item A Nash equilibrium is a strategy profile in which ...
  \begin{enumerate}
    \item no player ever randomises
    \item every player's strategy is a best response to the others' strategies
    \item every player plays a strictly dominant strategy
    \item the sum of the players' payoffs is maximised
  \end{enumerate}
  \item A strategy \(s^*\) for player \(i\) is a best response to the others if ...
  \begin{enumerate}
    \item it gives every player the same payoff
    \item it is the strategy with the largest support
    \item it gives player \(i\) at least as high a payoff as any other strategy, given the others
    \item it strictly dominates every other strategy
  \end{enumerate}
  \item In Rock-Paper-Scissors, the unique Nash equilibrium is for each player to ...
  \begin{enumerate}
    \item play each action with probability \(1/3\)
    \item copy the opponent's last move
    \item always play Rock
    \item there is no Nash equilibrium
  \end{enumerate}
  \item The support of a mixed strategy is ...
  \begin{enumerate}
    \item the set of actions played with positive probability
    \item the action with the highest payoff
    \item the opponent's best response
    \item the set of all of the player's actions
  \end{enumerate}
  \item An action is strictly dominated if ...
  \begin{enumerate}
    \item it is never a best response in a pure equilibrium
    \item the opponent also has a dominated action
    \item it gives a lower payoff than some other action against at least one opponent strategy
    \item some strategy gives a strictly higher payoff against every strategy of the opponents
  \end{enumerate}
  \item By the best response condition, a mixed strategy is a best response if and only if ...
  \begin{enumerate}
    \item the opponent is also indifferent
    \item all actions are played with equal probability
    \item it places all weight on a single action
    \item every action in its support yields the maximal expected payoff
  \end{enumerate}
  \item A strictly dominant action for a player is one that ...
  \begin{enumerate}
    \item is a best response to at least one opponent strategy
    \item maximises the combined payoff of both players
    \item gives a strictly higher payoff than every other action, against every opponent profile
    \item is played in some Nash equilibrium
  \end{enumerate}
  \item A pure strategy is a special case of a mixed strategy in which ...
  \begin{enumerate}
    \item all the probability is placed on a single action
    \item the player has only one action
    \item the opponent is indifferent
    \item every action is equally likely
  \end{enumerate}
  \item Nash's theorem guarantees that every finite game has ...
  \begin{enumerate}
    \item no Nash equilibrium unless there is a dominant strategy
    \item a Nash equilibrium only if it is zero sum
    \item a unique Nash equilibrium
    \item at least one Nash equilibrium, possibly in mixed strategies
  \end{enumerate}
  \item A coordination game (such as choosing the same side of the road) typically has ...
  \begin{enumerate}
    \item infinitely many pure equilibria
    \item no Nash equilibrium
    \item a single dominant-strategy equilibrium
    \item two pure equilibria and one mixed equilibrium
  \end{enumerate}
  \item To find the mixed equilibrium of a 2x2 game you ...
  \begin{enumerate}
    \item eliminate the dominant strategy
    \item maximise each player's own payoff
    \item set both players to play uniformly
    \item choose each player's mix so the other player is indifferent between their actions
  \end{enumerate}
  \item A two-player game is non-degenerate if ...
  \begin{enumerate}
    \item the payoff matrices sum to zero
    \item no strategy of support size \(k\) has more than \(k\) best response actions
    \item every player has a strictly dominant strategy
    \item it has a unique Nash equilibrium
  \end{enumerate}
  \item An action that is weakly (but not strictly) dominated ...
  \begin{enumerate}
    \item is never a best response
    \item can still be a best response against some opponent strategy
    \item is always played in equilibrium
    \item must be strictly dominated too
  \end{enumerate}
  \item The support enumeration algorithm finds Nash equilibria by ...
  \begin{enumerate}
    \item removing all dominated actions
    \item checking each possible pair of supports for the indifference and best response conditions
    \item simulating the game many times
    \item computing the value of the game
  \end{enumerate}
  \item If a game has a strictly dominant strategy for each player, then ...
  \begin{enumerate}
    \item there is no Nash equilibrium
    \item the equilibrium must be mixed
    \item the profile of dominant strategies is the unique Nash equilibrium
    \item there are many Nash equilibria
  \end{enumerate}
  \item The existence of a Nash equilibrium in every finite game is proved using ...
  \begin{enumerate}
    \item the central limit theorem
    \item the fundamental theorem of algebra
    \item Kakutani's fixed-point theorem applied to the best response correspondence
    \item linear programming duality
  \end{enumerate}
  \item The map \(\beta_i(s_{-i}) = \arg\max_{\sigma_i} u_i(\sigma_i, s_{-i})\) is called the ...
  \begin{enumerate}
    \item payoff matrix
    \item best response correspondence
    \item replicator equation
    \item support
  \end{enumerate}
  \item A mixed strategy for a player is ...
  \begin{enumerate}
    \item a probability distribution over that player's actions
    \item a payoff value
    \item an ordering of the opponents
    \item a single chosen action
  \end{enumerate}
  \item In the support enumeration algorithm, for a candidate pair of supports one solves equations making each player ...
  \begin{enumerate}
    \item prefer their first action
    \item minimise the opponent's payoff
    \item indifferent between the actions in their own support
    \item play a pure strategy
  \end{enumerate}
  \item Nash's theorem guarantees at least one equilibrium but ...
  \begin{enumerate}
    \item never a mixed one
    \item always exactly one
    \item only in zero-sum games
    \item not necessarily a unique one
  \end{enumerate}
\end{enumerate}
\end{document}
